\section{Special-function and kernel calculations}\label{app:special}
This appendix supplies the spectral calculation for Lemma~\ref{lem:fc}, the Gutzmer comparison used for the example
amplitudes, and the kernel bounds used in the cusp sharpness argument. The main estimates and their applications
are stated in the body of the paper.

\subsection{The first-order Cauchy estimate}\label{app:fc-proof}
\begin{proof}[Proof of Lemma~\ref{lem:fc}]
\medskip\noindent\textbf{Reduction.}

We first Fourier-transform in $x$. On the mode $\ue^{ikx}$, $k\ne0$, put $z=|k|y$ and write the function of $y$
as $z^{1/2}f(z)$, with $f\in L^2(dz/z)$. Then $-\Delta-\tfrac14=-(z\partial_z)^2+z^2$,
$X=i\,\operatorname{sgn}(k)\,z$ is multiplication, and $Y=z\partial_z+\tfrac12$.
Thus $Y_0:=Y-\tfrac12=z\partial_z=:D_z$ is skew-adjoint on $L^2(dz/z)$.

In the unitary Kontorovich--Lebedev picture
\[G=w(\nu)^{1/2}F,\qquad F(\nu)=\int_0^\infty K_{i\nu}(z)f(z)\,dz/z,\qquad w(\nu)=\tfrac{2}{\pi^2}\nu\sinh\pi\nu\]
(\cite{Y}), $\Lambda$ acts as multiplication by $\Lambda(\nu)=(\nu^2+\tfrac14)^{1/2}$.

The operators are the same for every $k$, so bounds are uniform in $k$. On the $k=0$ mode, $X=0$, while $Y_0$ and $\Lambda$ commute: in $r=\log y$, $Y_0=\partial_r$ and $\Lambda=(-\partial_r^2+\tfrac14)^{1/2}$. So
\[\norm{\ue^{\alpha''\Lambda}Y_0\ue^{-\alpha'\Lambda}}\le\norm{\Lambda\ue^{-(\alpha'-\alpha'')\Lambda}}\]
there. The mode $k=0$ has measure zero in the direct integral.

Put $\gamma:=\alpha'-\alpha''$ and write
\[
\ue^{\alpha''\Lambda}Z\ue^{-\alpha'\Lambda}=\bigl(\ue^{\alpha''\Lambda}Z\ue^{-\alpha''\Lambda}\Lambda^{-1}\bigr)\bigl(\Lambda\ue^{-\gamma\Lambda}\bigr).
\]
The second factor has norm $\le\sup_{\lambda\ge1/2}\lambda\ue^{-\gamma\lambda}\le1/(e\gamma)$. It remains to bound the first factor uniformly for $\alpha''\le\tau_0$, and $Y=Y_0+\tfrac12$ adds at most $\tfrac12$.

\medskip\noindent\textbf{Spectral notation.} Put $p:=\pi(\nu+\nu')/2$, $q:=\pi(\nu-\nu')/2$, $x:=\nu-\nu'$. Since $\sinh\pi\nu\,\sinh\pi\nu'=\sinh(p+q)\sinh(p-q)$,
\[
(w(\nu)w(\nu'))^{1/2}=\tfrac{2}{\pi^2}(\nu\nu')^{1/2}(\sinh^2p-\sinh^2q)^{1/2}=\tfrac{2}{\pi^2}(\nu\nu')^{1/2}(\cosh^2p-\cosh^2q)^{1/2}.
\]

\medskip\noindent\textbf{Mellin product formula.} For $\operatorname{Re}s>|\operatorname{Re}a|+|\operatorname{Re}b|$,
\[
M_s(a,b):=\int_0^\infty z^{s-1}K_a(z)K_b(z)\,dz=2^{s-3}\,\frac{\Gamma\bigl(\tfrac{s+a+b}{2}\bigr)\Gamma\bigl(\tfrac{s+a-b}{2}\bigr)\Gamma\bigl(\tfrac{s-a+b}{2}\bigr)\Gamma\bigl(\tfrac{s-a-b}{2}\bigr)}{\Gamma(s)}.
\]
\emph{Proof of the formula.} Write $K_b=K_{-b}$ (\cite[10.27.3]{DLMF}) and use the integral representation
\cite[10.32.10]{DLMF} for both Bessel factors. The integration in $z$ is
\[
\int z^{m-1}\ue^{-z^2c/4}\,dz=\tfrac12(4/c)^{m/2}\Gamma(m/2),
\qquad m=s+a-b,\quad c=1/u+1/v.
\]
Substitute $u=rt$ and $v=r(1-t)$ in the remaining double integral. The $r$-integral gives
$\Gamma(\tfrac{s-a+b}{2})$, and the $t$-integral gives the beta function
$\Gamma(\tfrac{s-a-b}{2})\Gamma(\tfrac{s+a+b}{2})/\Gamma(s)$, yielding the formula above.
The absolute integrand is the same expression with $a,b$ replaced by $\operatorname{Re}a,\operatorname{Re}b$ (since $|(z/2)^a|=(z/2)^{\operatorname{Re}a}$). For $\operatorname{Re}s>|\operatorname{Re}a|+|\operatorname{Re}b|$ the same computation then gives a finite value with all Gamma arguments of positive real part, so Tonelli justifies every exchange. This covers $a=i\nu$, $b=i\nu'$ at $s=1$, and $b=i\nu'\pm1$ at $s=1+\eta$, $\eta>0$. \qed

\medskip\noindent\textbf{(X) Kernel of $z$.}

Since the measure is $dz/z$, the kernel of multiplication by $z$ is $(w w')^{1/2}M_1(i\nu,i\nu')$.

The identity $|\Gamma(\tfrac12+iy)|^2=\pi/\cosh\pi y$ gives
\[
Zt_1(\nu,\nu')=(ww')^{1/2}\frac{\pi^2}{4\cosh p\cosh q}=\frac{(\nu\nu')^{1/2}(\cosh^2p-\cosh^2q)^{1/2}}{2\cosh p\cosh q},
\]
so $|Zt_1|\le(\nu\nu')^{1/2}/(2\cosh q)$.

The kernel of $\ue^{\alpha''\Lambda}z\ue^{-\alpha''\Lambda}\Lambda^{-1}$ is bounded by $\ue^{\alpha''|x|}(\nu\nu')^{1/2}/(2\cosh(\pi x/2)\Lambda(\nu'))$.

Since $(\nu\nu')^{1/2}\le(\nu+\nu')/2\le\nu'+|x|/2$ and $\Lambda(\nu')\ge\max(\nu',1/2)$, we have
\[
(\nu\nu')^{1/2}/\Lambda(\nu')\le1+|x|.
\]

So the kernel is at most $g(x):=(1+|x|)\ue^{\tau_0|x|}/(2\cosh(\pi x/2))$, which is integrable for $\tau_0<\pi/2$. The Schur test gives the bound $\int g$.

\medskip\noindent\textbf{(Y) Kernel of $D_z$.} The integral $\int K_{i\nu}K'_{i\nu'}\,dz$ does not converge at $z=0$, so we regularize.

Let $F,G\in C_c^\infty((0,\infty))$ be the KL transforms of $f,g$ (the core $\mathcal C$).

On $\mathcal C$, $f(e^{-L})$ and $zg'(z)|_{z=e^{-L}}$ decay faster than any power of $L$ as $L\to\infty$: they are integrals of smooth compactly supported densities against $\ue^{\pm i\nu L}$-oscillations, plus $O(\ue^{-2L})$. This uses $K_{i\nu}=\tfrac12\pi(I_{-i\nu}-I_{i\nu})/\sin(i\nu\pi)$ (\cite[10.27.4]{DLMF}) with the series of $I_{\pm i\nu}$ (\cite[10.25.2]{DLMF}), and $1/\sin(i\nu\pi)$ is smooth on $\operatorname{supp}F\subset(0,\infty)$.

Hence $\langle f,D_zg\rangle=\int_0^\infty\bar f g'\,dz=\lim_{\eta\to0^+}\int z^\eta\bar fg'\,dz$, by dominated convergence.

For $\eta>0$ Fubini applies, and $K'_\nu=-(K_{\nu-1}+K_{\nu+1})/2$ (\cite[10.29.1]{DLMF}, with $\mathcal Z_\nu=\ue^{\nu\pi i}K_\nu$ in DLMF's convention, under which 10.29.1 reads $\mathcal Z_{\nu-1}+\mathcal Z_{\nu+1}=2\mathcal Z'_\nu$ and gives the minus sign for $K$) gives
\[
\begin{gathered}
\langle f,D_zg\rangle=\lim_{\eta\to0^+}\iint\bar F(\nu)G(\nu')\,w(\nu)w(\nu')\,I_\eta(\nu,\nu')\,d\nu\,d\nu',\\
I_\eta=-\tfrac12\bigl[M_{1+\eta}(i\nu,i\nu'+1)+M_{1+\eta}(i\nu,i\nu'-1)\bigr].
\end{gathered}
\]
Both terms are admissible, since $\operatorname{Re}s=1+\eta>0+1$.

In $M_{1+\eta}(i\nu,i\nu'+1)$ the factor $\Gamma((\eta+ix)/2)=(2/(\eta+ix))\Gamma(1+(\eta+ix)/2)$ is singular, and the other factors are smooth. In $M_{1+\eta}(i\nu,i\nu'-1)$ the singular factor is $\Gamma((\eta-ix)/2)$.

So $I_\eta=\frac{2}{\eta+ix}A_\eta+\frac{2}{\eta-ix}B_\eta$, where $A_\eta,B_\eta$ are smooth in $(\nu,\nu',\eta)$ on the support.

At $\eta=0$ and $x=0$, the reflection identity $\Gamma(1+w)\Gamma(-w)=-\pi/\sin\pi w$ gives
\[
\begin{aligned}
A_0&=-\tfrac12\cdot\tfrac14\,\Gamma(1+ip/\pi)\Gamma(-ip/\pi)=-\tfrac18\cdot\tfrac{i\pi}{\sinh p},\\
B_0&=-\tfrac12\cdot\tfrac14\,\Gamma(ip/\pi)\Gamma(1-ip/\pi)=-\tfrac18\cdot\tfrac{-i\pi}{\sinh p}.
\end{aligned}
\]
So $A_0+B_0=0$ on the diagonal.

The Sokhotski--Plemelj formula gives $\frac{2}{\eta\pm ix}\to\mp2i\,\mathrm{PV}\frac1x+2\pi\delta(x)$. The $\delta$-terms carry $A_0+B_0$ at $x=0$, so they cancel.

For $x\neq0$ the pointwise limit is $I_0=-\pi^2/(4\sinh p\,\sinh q)$ (same Gamma reflections). Hence, on $\mathcal C\times\mathcal C$,
\[
\langle f,D_zg\rangle=\mathrm{PV}\iint\bar F G\,w\,w'\,I_0,
\qquad\text{i.e.}\qquad
\begin{gathered}
\tilde D(\nu,\nu')=-\frac{(\nu\nu')^{1/2}\rho(\nu,\nu')}{2\sinh q},\\
\rho:=\Bigl(1-\frac{\sinh^2q}{\sinh^2p}\Bigr)^{1/2}\in[0,1].
\end{gathered}
\]

This kernel is real and antisymmetric, as $D_z$ must be.

As a consistency check, not used in the proof, $(\nu^2-\nu'^2)\tilde D=-2Zt_2$, where
$Zt_2=(w(\nu)w(\nu'))^{1/2}M_2(\nu,\nu')$ is the kernel of $z^2$, with $M_2$ as in
Remark~\ref{rem:fc-y-alt}. This agrees with $[\Lambda^2,D_z]=[z^2,D_z]=-2z^2$.

\medskip\noindent\textbf{(Y) bound.} The kernel of $Q:=\ue^{\alpha''\Lambda}D_z\ue^{-\alpha''\Lambda}\Lambda^{-1}$ is $\ue^{\alpha''(\Lambda(\nu)-\Lambda(\nu'))}\tilde D(\nu,\nu')/\Lambda(\nu')$. It is valid on $\mathcal C$, because $\ue^{\pm\alpha''\Lambda}$ and $\Lambda^{-1}$ preserve $\mathcal C$. Split $Q=Q_{\rm far}+Q_{\rm near}$, with $Q_{\rm near}:=Q\cdot\mathbf 1_{|x|\le1}$.

\textbf{$Q_{\rm far}$.} Here $|Q|\le(1+|x|)\ue^{\tau_0|x|}/(2\sinh(\pi|x|/2))$ on $|x|>1$. This is integrable, so the Schur test applies.

\textbf{$Q_{\rm near}$.} Write $1/(2\sinh(\pi x/2))=1/(\pi x)+r(x)$, with $r$ bounded on $|x|\le1$.

The $r$-part has a bounded kernel supported in $|x|\le1$, so the Schur test applies.

In the main part $-a(\nu,\nu')/(\pi(\nu-\nu'))$, put
\[a(\nu,\nu'):=\ue^{\alpha''(\Lambda(\nu)-\Lambda(\nu'))}(\nu\nu')^{1/2}\rho(\nu,\nu')/\Lambda(\nu')\]
Then $a(\nu,\nu')=a(\nu',\nu')+[a(\nu,\nu')-a(\nu',\nu')]$, with $a(\nu',\nu')=\nu'/\Lambda(\nu')\in[0,1]$.

The diagonal part is the truncated Hilbert transform $\mathrm{PV}\,\frac{1}{\pi x}\mathbf 1_{|x|\le1}$
composed with multiplication by $\nu'/\Lambda(\nu')$. This transform is bounded on $L^2(\mathbb R)$:
its Fourier multiplier $-\tfrac{2i}{\pi}\operatorname{Si}(|\xi|)\operatorname{sgn}\xi$ has modulus
at most $(2/\pi)\operatorname{Si}(\pi)<1.2$. Extending functions by zero from $(0,\infty)$ therefore gives the required bound.

For the off-diagonal remainder, write $s(\nu):=(\nu\nu')^{1/2}/\Lambda(\nu')$ and
$E:=\ue^{\alpha''(\Lambda(\nu)-\Lambda(\nu'))}$. Then
\[
a(\nu,\nu')-a(\nu',\nu')=(E-1)s\rho+(\rho-1)s+(s(\nu)-s(\nu')).
\]
For $|x|\le1$, the first term is controlled by $|E-1|\le\tau_0\ue^{\tau_0}|x|$ and
$s\rho\le s\le1+|x|\le2$. For the second, $\sinh|q|/\sinh p\le|q|/p$ when $|q|\le p$ gives
$1-\rho\le\sinh^2q/\sinh^2p\le|q|/p=|x|/(\nu+\nu')$, and hence
$(1-\rho)s\le|x|(\nu\nu')^{1/2}/((\nu+\nu')\Lambda(\nu'))\le|x|$.
Finally, $|s(\nu)-s(\nu')|=|\sqrt\nu-\sqrt{\nu'}|\sqrt{\nu'}/\Lambda(\nu')\le|x|^{1/2}$,
since $\sqrt{\nu'}/\Lambda(\nu')\le1$.
So the remainder kernel is $\le C|x|^{-1/2}$ on $|x|\le1$, which is integrable. The Schur test therefore bounds the remainder.
Combining the near and far bounds gives $\norm{Q}\le C(\tau_0)$ on $\mathcal C$, so $Q$ extends to a bounded operator. The bound for $\ue^{\alpha''\Lambda}D_z\ue^{-\alpha'\Lambda}$ follows on the dense set $\Lambda\ue^{-\gamma\Lambda}\mathcal C$, and then by closure. \qedhere
\end{proof}

\begin{remark}[Proof of (Y) without a principal value]\label{rem:fc-y-alt}
The case $Z=Y$ of Lemma~\ref{lem:fc} also follows without the kernel of $D_z$. For $A,B\in C_c^\infty((0,\infty))$ put $a=wA$, $b=wB$, $\hat a(t)=\int a(\nu)\ue^{i\nu t}\,d\nu$, $\psi_a(z)=\int a(\nu)K_{i\nu}(z)\,d\nu=\tfrac12\int\ue^{-z\cosh t}\hat a(t)\,dt$ and $S(A,B)=\langle\psi_a,D_z\psi_b\rangle_{L^2(dz/z)}$.
\begin{itemize}
\item By Fubini, $S(A,B)=-\tfrac14\iint\overline{\hat a(t)}\,\hat b(u)\,\frac{\cosh u}{\cosh t+\cosh u}\,dt\,du$, absolutely convergent.
\item Since $(\partial_u^2-\partial_t^2)\frac{\cosh u}{\cosh t+\cosh u}=2(\cosh t+\cosh u)^{-2}$, integrating by parts twice in each variable gives
\[S(\Lambda^2A,B)-S(A,\Lambda^2B)=-2\iint\bar A(\nu)B(\nu')\,w(\nu)w(\nu')\,M_2(\nu,\nu')\,d\nu\,d\nu',\]
where $M_2(\nu,\nu')=\int_0^\infty zK_{i\nu}K_{i\nu'}\,dz=\tfrac12\,\mathsf g(\tfrac{\nu+\nu'}2)\,\mathsf g(\tfrac{\nu-\nu'}2)$ and $\mathsf g(y)=\pi y/\sinh\pi y$ (the Mellin product formula at $s=2$).
\item The multipliers $\ue^{\pm s\Lambda}$ preserve $C_c^\infty((0,\infty))$. Integrating first along $\ue^{-t\Lambda}$ in both arguments over $t\in(0,\infty)$, then along $(\ue^{s\Lambda},\ue^{-s\Lambda})$ over $s\in(0,a)$, gives
\[S(\ue^{a\Lambda}A,\ue^{-a\Lambda}B)-S(A,B)=\iint\bar A(\nu)B(\nu')\,w(\nu)w(\nu')\,R_a(\nu,\nu')\,d\nu\,d\nu',\]
with $R_a=-2M_2\,\phi_a(\Lambda(\nu)-\Lambda(\nu'))/(\Lambda(\nu)+\Lambda(\nu'))$ and $\phi_a(\delta)=\int_0^a\ue^{s\delta}\,ds$.
\item Hence, on products in the KL picture and then by density, $\ue^{a\Lambda}Y\ue^{-(a+\gamma)\Lambda}=Y\ue^{-\gamma\Lambda}+T_{R_a}\ue^{-\gamma\Lambda}$, where $T_RG(\nu)=\int R(\nu,\nu')G(\nu')w(\nu')\,d\nu'$. The first term has norm at most $2/\gamma$, by $\norm{Yu}^2\le\norm{u}\,\norm{\Delta u}$. For the second, $\mathsf g(p')\cosh(\pi p')\le1+\pi p'$ with $p'=(\nu+\nu')/2$, together with the inequalities used in (X), gives
\[(w(\nu)w(\nu'))^{1/2}\,|R_a(\nu,\nu')|\,\ue^{-\gamma\Lambda(\nu')}\le C(\tau_0)\,\gamma^{-1}(1+|x|)\ue^{\tau_0|x|}\,\mathsf g(x/2),\qquad x=\nu-\nu',\]
which is integrable for $\tau_0<\pi/2$, and the Schur test concludes.
\end{itemize}
\end{remark}

\subsection{A comparison by the Gutzmer identity}\label{subsec:gutzmer}
This subsection is used only for Lemma~\ref{lem:ex}. Let $G=SL(2,\mathbb R)$, $h_0=\mathrm{diag}(1,-1)$ and $x_0=i$, so that
$d(x_0,\exp(th_0)x_0)=2|t|$. For $0<R\le\pi/4$ let $\Xi_R=G\exp(i(-R,R)h_0)x_0$, a $G$-invariant domain in the complexification
$G_{\mathbb C}/K_{\mathbb C}$ of $\Hb=G/K$; $\Xi_{\pi/4}$ is the crown \cite[\S7.2]{K}. For a holomorphic function $f$ on $\Xi_R$ and $0\le r<R$ put
\[
N_r(f)^2:=\int_G|f(g\exp(irh_0)x_0)|^2\,dg ,
\]
with the Haar measure normalized so that $N_0(f)=\norm f_{L^2(\Hb)}$.

\begin{lemma}[comparison with the orbit norm]\label{lem:gutzmer}
Let $0<r<R\le\pi/4$ with $r\le\tfrac12$.
\begin{itemize}
\item[(a)] Let $f\in L^2(\Hb)$ extend holomorphically to $\Xi_R$ with $\sup_{0\le\rho<R}N_\rho(f)<\infty$. Then
$f\in\operatorname{Dom}(\ue^{r\Lambda})$ and $\norm{\ue^{r\Lambda}f}\le(\pi\ue)^{1/2}N_r(f)$.
\item[(b)] The function $X_0=\cosh d(x_0,\cdot)$ extends holomorphically to $\Xi_{\pi/4}$, and for $g\in G$ its value at
$g\exp(irh_0)x_0$ is $\cos(2r)\,A+i\sin(2r)\,B$, where $A=X_0(gx_0)\ge1$ and $B$ is real with $B^2\le A^2-1$.
\end{itemize}
\end{lemma}
\begin{proof}
(b) Realize $\Hb$ as the sheet $p_0>0$ of the hyperboloid $\{p\in\mathbb R^3:-p_0^2+p_1^2+p_2^2=-1\}$, with $x_0=e_0=(1,0,0)$, the group $G$
acting through $SO_0(2,1)$, and $G_{\mathbb C}/K_{\mathbb C}$ as the complex quadric given by the same equation. Then $X_0(p)=p_0$ is a linear function,
$\exp(th_0)e_0=(\cosh2t,\sinh2t,0)$, and so $\exp(irh_0)e_0=\cos(2r)\,e_0+i\sin(2r)\,e_1$. Hence
$X_0(g\exp(irh_0)x_0)=\cos(2r)\,(ge_0)_0+i\sin(2r)\,(ge_1)_0$. Here $(ge_0)_0=X_0(gx_0)=A$, and $q:=ge_1$ is a real vector with
$-q_0^2+q_1^2+q_2^2=1$ that is orthogonal to $p:=ge_0$ for the form: $Aq_0=p_1q_1+p_2q_2$. By the Cauchy--Schwarz inequality
$A^2q_0^2\le(A^2-1)(1+q_0^2)$, that is $q_0^2\le A^2-1$.

(a) In the spectral decomposition of $L^2(\Hb)$ by the Helgason--Fourier transform, $-\Delta$ acts as $\lambda^2+\tfrac14$, $\lambda\ge0$, and
$\Lambda$ as $(\lambda^2+\tfrac14)^{1/2}\le\lambda+\tfrac12$. The Gutzmer identity (\cite[\S7.2, the identity before Theorem~7.2]{K}, where the
spectral parameter is $2\lambda$) states that for $f$ as in (a) and $r<R$ the number $N_r(f)^2$ is the integral of $|\hat f|^2$ against the
Plancherel measure with the additional weight $\phi_\lambda(\exp(2irh_0))$, the value at $\exp(2irh_0)x_0$ of the holomorphic extension of the
spherical function $\phi_\lambda$. A radial eigenfunction of $\Delta$ with eigenvalue $-(\lambda^2+\tfrac14)$ that equals $1$ at $x_0$ is the regular
solution of Legendre's equation in the variable $X_0$, so $\phi_\lambda=P_{-1/2+i\lambda}(X_0)$, a conical function \cite[\S14.20]{DLMF}; by (b),
$X_0(\exp(2irh_0)x_0)=\cos4r$. With $\theta:=4r\in(0,2]$, the Mehler--Dirichlet integral \cite[14.20.9]{DLMF} and the substitution $\phi=\theta v$ give
\[
P_{-1/2+i\lambda}(\cos\theta)=\frac2\pi\int_0^1\frac{\cosh(\lambda\theta v)\,\theta}{\sqrt{2(\cos\theta v-\cos\theta)}}\,dv
\ge\frac2\pi\int_0^1\frac{\cosh(\lambda\theta v)}{\sqrt{1-v^2}}\,dv=I_0(\lambda\theta),
\]
because $2(\cos\theta v-\cos\theta)=4\sin\tfrac{\theta(1+v)}2\sin\tfrac{\theta(1-v)}2\le\theta^2(1-v^2)$; the last equality is
\cite[10.32.1]{DLMF}. From the same formula, $I_0(x)=\tfrac1\pi\int_0^\pi\ue^{x\cos t}dt\ge\tfrac1\pi\int_0^{(1+x)^{-1/2}}\ue^{x(1-t^2/2)}dt
\ge(\pi\sqrt{\ue})^{-1}(1+x)^{-1/2}\ue^{x}$ for $x\ge0$. Since $\ue^{4r\lambda}\ge(1+4r\lambda)^{1/2}\ue^{2r\lambda}$ and
$\ue^{2r\Lambda}\le\ue^{r}\ue^{2r\lambda}\le\sqrt{\ue}\,\ue^{2r\lambda}$,
\[
\phi_\lambda(\exp(2irh_0))\ge(\pi\sqrt{\ue})^{-1}\ue^{2r\lambda}\ge(\pi\ue)^{-1}\ue^{2r(\lambda^2+1/4)^{1/2}} ,
\]
and integrating against $|\hat f|^2$ gives (a).
\end{proof}

\subsection{Kernels of the Poisson and heat semigroups}

The two lemmas of this subsection are used in Section~\ref{subsec:cl} (Lemma~\ref{cl:kernel}). They are not used in the proof of
Theorem~\ref{pc:main}.
Let $\mu_\tau(s)=\tau(4\pi)^{-1/2}s^{-3/2}e^{-\tau^2/(4s)}$ (so that $e^{-\tau\osLam}=\int_0^\infty\mu_\tau(s)e^{s\Delta}\,ds$, by
the subordination identity $e^{-\tau\sqrt A}=\int_0^\infty\mu_\tau(s)e^{-sA}ds$, $A\ge0$), and let $p_s$ be the heat kernel of $\Hb$, which is
given by McKean's formula
\[
p_s(d)=\frac{\sqrt2\,e^{-s/4}}{(4\pi s)^{3/2}}\int_d^\infty\frac{\rho\,e^{-\rho^2/(4s)}}{(\cosh\rho-\cosh d)^{1/2}}\,d\rho
\]
(\cite[(1.7)]{GN98}). Let
$k_\tau(r)=\int_0^\infty\mu_\tau(s)p_s(r)\,ds$ be the radial kernel of $e^{-\tau\osLam}$.

\begin{lemma}[Poisson kernel representation and bounds]\label{os:poisson}
For $0<\tau\le\tau_0<\pi/2$ and $r\ge0$,
\[
k_\tau(r)=\frac{\sqrt2}{8\pi^2}\,\tau\int_r^\infty\frac{\rho\,K_2\bigl(\sqrt{\rho^2+\tau^2}/2\bigr)}{\rho^2+\tau^2}\,\frac{d\rho}{(\cosh\rho-\cosh r)^{1/2}}.
\]
Consequently there are $0<c_\tau,C_\tau<\infty$, depending only on $\tau_0$, with
\begin{align}
k_\tau(r)&\le C\,\tau\,e^{-r}\osbr{r}^{-3/2}\quad(r\ge1),\label{os:ub}\\
k_\tau(r)&\ge c\,\tau\,e^{-r}\osbr{r}^{-3/2}\quad(r\ge1),\label{os:lb}\\
c\,\tau(\tau^2+r^2)^{-3/2}\le k_\tau(r)&\le C\,\tau(\tau^2+r^2)^{-3/2}\quad(0\le r\le1).\label{os:small}
\end{align}
\end{lemma}
\begin{proof}
Insert the formula for $p_s$ into $k_\tau=\int\mu_\tau p_s\,ds$ and interchange the integrals (Tonelli). The $s$-integral is
\[
\int_0^\infty s^{-3}e^{-(\rho^2+\tau^2)/(4s)-s/4}\,ds=2\frac{1/4}{(\rho^2+\tau^2)/4}K_2\Bigl(\frac{\sqrt{\rho^2+\tau^2}}2\Bigr)
=\frac{2K_2(\sqrt{\rho^2+\tau^2}/2)}{\rho^2+\tau^2},
\]
by $\int_0^\infty s^{\nu-1}e^{-\beta s-\gamma/s}ds=2(\gamma/\beta)^{\nu/2}K_\nu(2\sqrt{\beta\gamma})$ (\cite{DLMF}, 10.32.10, and the
standard Bessel integral) with $\nu=-2$, $\beta=\tfrac14$, $\gamma=(\rho^2+\tau^2)/4$, and $K_{-2}=K_2$. The constant is
$\tfrac{\tau}{2\sqrt\pi}\cdot\sqrt2\cdot(4\pi)^{-3/2}\cdot2=\tfrac{\sqrt2}{8\pi^2}\tau$. This gives the representation.

For the bounds write $\rho=r+v$, so $\cosh\rho-\cosh r=2\sinh(r+v/2)\sinh(v/2)$. For $r\ge1$ the factor
$K_2(\sqrt{\rho^2+\tau^2}/2)$ is comparable to $e^{-\rho/2}\rho^{-1/2}$ uniformly for $\tau\le\tau_0$, the factor $\rho/(\rho^2+\tau^2)$
is at most $\rho^{-1}\le r^{-1}$, and for $0\le v\le1$ it is at least $((1+\tau_0^2)\,2r)^{-1}$; also $(\cosh\rho-\cosh r)^{-1/2}=(2\sinh(r+v/2)\sinh(v/2))^{-1/2}$ is at most $Ce^{-r/2}(\sinh(v/2))^{-1/2}$ for all $v>0$ and comparable to it for $0<v\le1$. The $v$-integral of
$e^{-v/2}(\sinh(v/2))^{-1/2}$ over $(0,\infty)$ is finite, which gives \eqref{os:ub}; its integral over $(0,1]$ is positive, and restricting to $0\le v\le1$ gives \eqref{os:lb}. For $r\le1$ split the
integral at $\rho=2$; there $\cosh\rho-\cosh r$ is comparable to $\rho^2-r^2$, and with $v=\sqrt{\rho^2-r^2}$ the main part is
comparable to $\int_0^{C}(r^2+\tau^2+v^2)^{-2}\,dv\asymp(r^2+\tau^2)^{-3/2}$, while the tail $\rho\ge2$ is bounded by the
exponential estimate. This gives \eqref{os:small}.
\end{proof}

Let $\osHK$ denote the continuation of $p_s$ to complex distance $\tilde d=a+i\theta$, $a\ge0$, $|\theta|\le\theta_0<\pi/2$.
McKean's formula for $p_s$ is first put in $d^2$-form; in the rest of this subsection $\tau$ denotes an integration variable, not the
parameter of Lemma~\ref{os:poisson}. With $g(w)=\sinh\sqrt w/(2\sqrt w)$ (an entire function of $w$) and
$F(w,\tau^2)=\int_0^1g(w+\lambda\tau^2)\,d\lambda=(\cosh\sqrt{w+\tau^2}-\cosh\sqrt w)/\tau^2$, substituting $\rho=\sqrt{d^2+\tau^2}$ gives
\[
p_s(d)=c_s\,e^{-d^2/(4s)}\int_0^\infty e^{-\tau^2/(4s)}F(d^2,\tau^2)^{-1/2}\,d\tau,\qquad c_s=\frac{\sqrt2e^{-s/4}}{(4\pi s)^{3/2}}.
\]

\begin{lemma}[Continued heat kernel]\label{os:hk}
Let $0<\theta_0<\pi/2$ and $c_{\theta_0}=\cos\theta_0\,e^{-\theta_0}$. For $a\ge0$, $|\theta|\le\theta_0$ and $s>0$,
\[
|\osHK_s(a+i\theta)|\le c_{\theta_0}^{-1/2}\,e^{\theta^2/(4s)}\,p_s(a).
\]
\end{lemma}
\begin{proof}
Put $d^2\mapsto\tilde d^2=a^2-\theta^2+2ia\theta$ in the $d^2$-representation. Let $x=a^2+\lambda\tau^2\ge a^2$ and $w=x-\theta^2+2ia\theta$,
and write $\sqrt w=A+iB$ with $A\ge0$. Then $|w|\le x+\theta^2$, so $B^2=(|w|-(x-\theta^2))/2\le\theta^2$; also $A^2=x-\theta^2+B^2\le x$.
Hence $0\le\sqrt x-A$, and if $x\ge\theta^2$ then $\sqrt x-A\le|\theta|$.

Now $\operatorname{Re}\frac{\sinh\zeta}{\zeta}=\frac{A\sinh A\cos B+B\cosh A\sin B}{A^2+B^2}$ for $\zeta=A+iB$. Since $|B|<\pi/2$ we have
$B\sin B\ge B^2\cos B$ (for either sign of $B$), and $\cosh A\ge\sinh A/A$, so $\operatorname{Re}(\sinh\zeta/\zeta)\ge\cos B\,\sinh A/A\ge\cos\theta_0\,\sinh A/A$. The
logarithmic derivative of $\sinh u/u$ lies in $[0,1]$, so $\sinh A/A\ge e^{-|\theta|}\sinh\sqrt x/\sqrt x$. Therefore
\[
\operatorname{Re}g(w)\ge c_{\theta_0}\,g(x)>0.
\]
(The case $x<\theta^2$ is handled directly, since then $\sinh A/A\ge1\ge(\sinh\sqrt x/\sqrt x)e^{-\theta_0}$.)

Hence $\operatorname{Re}F(\tilde d^2,\tau^2)\ge c_{\theta_0}F(a^2,\tau^2)>0$. For $\operatorname{Re}z>0$ one has $|z^{-1/2}|\le(\operatorname{Re}z)^{-1/2}$, so
$|F(\tilde d^2,\tau^2)^{-1/2}|\le c_{\theta_0}^{-1/2}F(a^2,\tau^2)^{-1/2}$. Finally
$|e^{-\tilde d^2/(4s)}|=e^{\theta^2/(4s)}e^{-a^2/(4s)}$, and the $\tau$-integral is dominated uniformly on compact sets of $\tilde d$, so the continuation is
holomorphic. This gives the bound.
\end{proof}
